\documentclass{article}
\usepackage{enumitem}

\begin{document}

\section*{Group 1: Morning Book Study (Alternate Focus Areas)}

You will alternate between the following books each morning, focusing on key topics and taking notes:

\begin{enumerate}
    \item \textbf{Valvano’s Book}:
    \begin{itemize}
        \item Read the first 40 pages (introductory concepts on computers and electronics).
        \item Ongoing study of embedded systems, focusing on how it ties into your practical projects.
    \end{itemize}

    \item \textbf{Yiu’s "Embedded Systems"}:
    \begin{itemize}
        \item Focus on chapters related to ADC and microcontroller-specific concepts.
        \item Supplement Valvano's book for hardware configuration and register-level understanding.
    \end{itemize}

    \item \textbf{"The Art of Electronics"}:
    \begin{itemize}
        \item Supplement Valvano with foundational electronics concepts, especially on signals, sinusoidal functions, and ADC.
        \item Focus on introductory chapters that explain electronics fundamentals, and reference page 14 on sinusoidal signals.
    \end{itemize}

    \item \textbf{Brown’s Book}:
    \begin{itemize}
        \item Use as a reference for minimal projects and automation, particularly ADC and hardware setup.
        \item Supplement Yiu and Valvano for specific topics like ADC and signal processing.
    \end{itemize}

    \item \textbf{CLRS and Sedgewick}:
    \begin{itemize}
        \item Continue with your focus on sliding window algorithms, KMP, Rabin-Karp, tries, and B-trees.
        \item De-prioritize the rest of the data structures since you're already familiar with them.
    \end{itemize}
\end{enumerate}

\section*{Group 2: Follow-Up on FFT Applied to ADC}

\begin{enumerate}
    \item \textbf{Review Polynomials and Coefficients in Relation to FFT}:
    \begin{itemize}
        \item \textbf{Reference Books}: CLRS and Sedgewick. Review sections on polynomials and their role in algorithms, focusing on their application in FFT.
        \item \textbf{Action}: Understand how polynomials form the basis for the FFT algorithm (e.g., how FFT decomposes a signal into sinusoidal components).
        \item \textbf{Goal}: Grasp the mathematical foundation (polynomials, coefficients) that underpins the FFT.
    \end{itemize}
    
    \item \textbf{Basic Electronics and ADC Fundamentals}:
    \begin{itemize}
        \item \textbf{Reference Books}: The Art of Electronics, Valvano’s Book, Yiu’s "Embedded Systems", and Brown’s Book.
        \item \textbf{Action}: Map out how the ADC works, particularly in the STM32 context. Extract formulas from the books and locate corresponding sections in the STM32 datasheets/manuals.
        \item \textbf{Goal}: Develop a clear understanding of ADC operations and their connection to signal processing with FFT.
    \end{itemize}

    \item \textbf{Sinusoidal Signals and Magnitude}:
    \begin{itemize}
        \item \textbf{Reference}: The Art of Electronics, page 14.
        \item \textbf{Action}: Learn the meaning of “magnitude” in the context of FFT and how it represents frequency components.
        \item \textbf{Goal}: Understand the breakdown of sinusoidal signals via FFT and the interpretation of magnitude.
    \end{itemize}

    \item \textbf{FFT Algorithm Based on Sedgewick}:
    \begin{itemize}
        \item \textbf{Reference}: Sedgewick’s Algorithms.
        \item \textbf{Action}: Continue using Sedgewick’s FFT implementation in your project. Verify the algorithm with real-time ADC data.
        \item \textbf{Goal}: Ensure the FFT algorithm works effectively in your project.
    \end{itemize}

    \item \textbf{Research the CMSIS FFT Library (for reference only)}:
    \begin{itemize}
        \item \textbf{Action}: Explore CMSIS FFT to familiarize yourself with the professional tools available for future reference.
        \item \textbf{Goal}: Gain awareness of how CMSIS libraries could be leveraged in the future.
    \end{itemize}

    \item \textbf{Formulas and Practical Applications}:
    \begin{itemize}
        \item \textbf{Action}: Gather key formulas related to ADC (e.g., converting ADC values to voltages, scaling signals for FFT input) and confirm their presence in the STM32 manual.
        \item \textbf{Goal}: Have practical formulas ready for your project.
    \end{itemize}
\end{enumerate}

\section*{Group 3: Algorithms and Data Structures}

\subsection*{1.1 Algorithms}
\begin{itemize}
    \item \textbf{Fast Fourier Transform (FFT)}: Study the algorithm for frequency domain analysis, focusing on efficient computation in real-time systems.
    \item \textbf{Sliding Window Algorithm}: Used for maintaining a subset of items (e.g., moving average or window-based data aggregation) in real-time data streams.
    \item \textbf{Pattern Matching}:
    \begin{itemize}
        \item KMP Algorithm: Knuth-Morris-Pratt for efficient string matching.
        \item Rabin-Karp Algorithm: Used for detecting patterns in streams of data.
    \end{itemize}
\end{itemize}

\subsection*{1.2 Data Structures}
\begin{itemize}
    \item \textbf{Trie}: Efficient for storing and querying dynamic sequences or strings.
    \item \textbf{B-tree}: Balanced tree structure used for large datasets.
\end{itemize}

\subsection*{2. Integrated Embedded Systems Tasks (Hardware + Algorithms)}
\subsubsection*{2.1 Priority 1: Foundation-Level Tasks}
\begin{itemize}
    \item \textbf{Temperature and Voltage Monitoring}
    \begin{itemize}
        \item Data Source: Internal temperature sensor, voltage reference (VREFINT)
        \item Data Structures: Ring Buffer, Priority Queue
        \item Algorithms: Moving Average, Exponential Smoothing, Threshold Detection
    \end{itemize}
    
    \item \textbf{Watchdog Timers}
    \begin{itemize}
        \item Data Source: Independent and Window Watchdog Timers (IWDG, WWDG)
        \item Data Structures: Circular Queue
        \item Algorithms: Timeout Management, Monitoring and Alerting
    \end{itemize}
    
    \item \textbf{Real-Time Clock (RTC) \& Timestamps}
    \begin{itemize}
        \item Data Source: RTC, SysTick timer
        \item Data Structures: Binary Search Tree, B-tree
        \item Algorithms: Scheduling Algorithms (Round-robin, priority-based), Binary Search
    \end{itemize}
\end{itemize}

\subsubsection*{2.2 Priority 2: Intermediate-Level Tasks}
\begin{itemize}
    \item \textbf{Interrupt-Driven Task Scheduling}
    \begin{itemize}
        \item Data Source: Timers, interrupts
        \item Data Structures: Priority Queue, Doubly Linked List
        \item Algorithms: Round-Robin Scheduling, Earliest Deadline First (EDF)
    \end{itemize}
    
    \item \textbf{Real-Time Signal Processing with ADC}
    \begin{itemize}
        \item Data Source: Analog-to-Digital Converter (ADC)
        \item Data Structures: Circular Queue, FFT Buffer
        \item Algorithms: FFT, Sliding Window, Finite State Machine (FSM) for pattern recognition
    \end{itemize}
    
    \item \textbf{Fault Monitoring and Recovery}
    \begin{itemize}
        \item Data Source: Non-Maskable Interrupt (NMI), system faults
        \item Data Structures: HashMap, Circular Queue
        \item Algorithms: Error Detection (Hamming Code), CRC
    \end{itemize}
\end{itemize}

\subsubsection*{2.3 Priority 3: Advanced-Level Tasks}
\begin{itemize}
    \item \textbf{DMA for Data Streaming}
    \begin{itemize}
        \item Data Source: DMA controller, ADC
        \item Data Structures: Ring Buffer, Trie
        \item Algorithms: Pattern Matching (KMP, Rabin-Karp), Circular Buffer Management
    \end{itemize}
    
    \item \textbf{Memory Management and Flash Programming}
    \begin{itemize}
        \item Data Source: Flash memory, EEPROM
        \item Data Structures: HashMap, Linked List
        \item Algorithms: Wear-Leveling, CRC for Data Integrity
    \end{itemize}
    
    \item \textbf{Power Consumption Monitoring}
    \begin{itemize}
        \item Data Source: Internal power sensors
        \item Data Structures: Priority Queue, HashMap
        \item Algorithms: Power Optimization (Dynamic Voltage/Frequency Scaling), Energy Profiling
    \end{itemize}
\end{itemize}

\end{document}
